% Round20 inertia20: an explicit failed induction, not a Green counterexample.
\section{Why higher products must retain the actual deletion model}
\label{r20he:section}

\paragraph{Claim and scope.}
The binary staircase proof does not extend by transporting every
surface transcendental projector only from the original unmodified
component product. Such a procedure can miss a nonzero exceptional
copy of a surface transcendental coefficient. The calculation below
is an obstruction to that proof procedure, not to strong G and not
to the more precise R9 deletion-model candidate.

Let $V$ be a smooth projective surface over a finite field $k$ with
$H^2_{\rm tr}(V,\mathbf Q_\ell)\ne0$. Let $D_1,D_2$ be smooth
projective surfaces and let $F_i\subset D_i$ be smooth projective
connected curves. Put
\[
 B=D_1\times D_2,\quad Z=F_1\times F_2,\quad
 b:W=\operatorname{Bl}_Z B\longrightarrow B,
 \quad f=1_V\times b:V\times W\longrightarrow V\times B.
 \tag{HE20.1}
\]
The codimension-two blowup formula is an actual Chow-motive
decomposition
\[
 h(V\times W)=h(V\times B)\oplus h(V\times Z)(1),
 \tag{HE20.2}
\]
where $(1)$ denotes the Lefschetz-motive factor, equivalently
the cohomological summand $H^{m-2}(V\times Z)(-1)$ in degree $m$.
Its exceptional injection and retraction are actual algebraic
correspondences. The first projector is $f^*f_*$, and $f_*$ is
zero on the exceptional summand.

Let $p_V$ be the surface transcendental cohomological projector.
The projector transported from the unmodified three-factor base is
\[
              P_{\rm base}=f^*(p_V\boxtimes1_B)f_*.
 \tag{HE20.3}
\]
It is zero on the entire second summand of (HE20.2). In particular
it misses the nonzero exceptional subspace
\[
 H^2_{\rm tr}(V)\otimes H^0(F_1)\otimes H^0(F_2)(-1)
                 \subset H^4(V\times W).
 \tag{HE20.4}
\]
The nonzero assertion follows from connectedness of the two curves,
the assumed nonzero transcendental space, and the directness of
the blowup decomposition.

By contrast the coordinate-deletion projector, which keeps the
\emph{actual} remaining component $W$, is
\[
                        P_{\rm del}=p_V\boxtimes1_W.
 \tag{HE20.5}
\]
It is the identity on (HE20.4). Thus
$P_{\rm base}\ne P_{\rm del}$ on full cohomology, even though
they agree over the locus where $b$ is an isomorphism. Equality
of two arbitrary projectors on a dense open cannot replace the
relative-support and coordinate-deletion calculation.

\paragraph{Realization in the actual semistable product geometry.}
The modification $b$ in (HE20.1) is exactly the component blowup
at the lower or upper diagonal vertex of the two-node product
chart: the center is the product of its two double curves.
For instance take two semistable curve models and form their
products with $\mathbf P^1_R$; their surface double strata
are copies of $\mathbf P^1_k$. The ordered binary product has
such a component blowup at every selected pair of curve nodes.
Disjoint further node centers merely add further direct summands.
Taking the product with a smooth projective model of $V$ gives
(HE20.2) as an actual component decomposition. For a concrete
nonzero coefficient, take over $\mathbf Z_5$ the smooth elliptic
curve $E:y^2=x^3+x+1$ and let $V=E_k\times E_k$. Counting its
residue points gives $\#E(\mathbf F_5)=9$, so its Frobenius
polynomial is $t^2+3t+5$ and the two roots have $5$-adic valuations
zero and one. The two normalized eigenvalues $\alpha^2/5$ and
$\beta^2/5$ in $H^2(V)(1)$ have nonzero $5$-adic valuations and
are not roots of unity. Every geometric divisor is defined over a
finite constant extension, so its normalized Frobenius eigenvalues
are roots of unity. These two eigenspaces therefore lie in the
nonzero transcendental summand. This example serves only to make
(HE20.4) visibly nonzero; its generic product is not claimed as a
new Green domain.

\paragraph{The exact remaining induction gate.}
For three or more bad surface factors, exceptional centers can
retain transcendental coefficients of untouched surface factors.
One must prove that the transported projectors from all actual
coordinate-deletion models capture those copies and that the
residual projector factors through abelian-type motives. Binary
abelianity of the new exceptional summand does not establish this
statement: in (HE20.2), the new summand contains the entire
$h(V)$ coefficient. No assertion that this particular coefficient
is not of abelian type is made. The missing issue is a controlled
factorization and incidence compatibility, not a claimed failure
of a standard conjecture.

The historical R9 GL/DG construction explicitly recognized the
distinction between (HE20.3) and (HE20.5). This note supplies a
concrete check of that distinction. It does not by itself validate
the general DG residual-core theorem or the uniform DT transition
theorem. Its only external algebraic inputs are the smooth blowup
formula and the surface transcendental projector already checked
in PX20; all maps above are displayed actual correspondences.
